\documentclass[11pt,a4paper]{article}

% ---------------------------------------------------------
% PREAMBLE (stable, warning-free, Overleaf-safe)
% ---------------------------------------------------------
\usepackage[utf8]{inputenc}
\usepackage[T1]{fontenc}
\usepackage{lmodern}
\usepackage{amsmath, amssymb, amsfonts}
\usepackage{bm}
\usepackage{geometry}
\usepackage{graphicx}
\usepackage{hyperref}
\usepackage{cite}
\usepackage{authblk}
\usepackage{physics}
\usepackage{mathtools}

\geometry{margin=2.4cm}

\title{\textbf{Companion XXXVI: The Baryonic Cycle in the Relaxation-Driven Cyclic Cosmology}}
\author{Michael Lehmann \\ \texttt{mi.lehmann@gmx.de}}
\date{April 2026}

\begin{document}
\maketitle

\begin{abstract}
The formation of galaxies is governed by the complex interplay between gas cooling, 
star formation, and feedback. These processes form a self-regulating baryonic cycle 
that determines the luminosity, morphology, and chemical evolution of galaxies. 
In the Relaxation-Driven Cyclic Cosmology (RDCC), the scale-dependent evolution of 
the gravitational potential modifies the thermodynamic pathways through which gas 
accretes, cools, and condenses into stars. This Companion develops a continuous 
description of the baryonic cycle in RDCC, deriving how the relaxation scale 
influences cooling times, feedback efficiency, and the equilibrium between inflows 
and outflows. The resulting framework provides a natural explanation for the 
suppressed star formation in low-mass halos and the early assembly of massive 
galaxies observed in deep surveys.
\end{abstract}
% ---------------------------------------------------------
\section{Introduction}

The baryonic cycle is the engine of galaxy formation. Gas flows into halos along 
filaments, cools radiatively, condenses into molecular clouds, forms stars, and is 
then reheated or expelled by stellar and AGN feedback. This cycle is not linear but 
self-regulating: the rate at which gas cools determines the star formation rate, 
which in turn determines the strength of feedback, which then modifies the cooling 
rate.

In RDCC, the gravitational potential that shapes this cycle evolves differently 
from the standard cosmological model. The relaxation mechanism suppresses the 
growth of small-scale modes and enhances the coherence of large-scale modes. This 
modifies the density, temperature, and entropy structure of the gas as it accretes 
onto halos. The baryonic cycle becomes sensitive to the relaxation scale 
$k_{\mathrm{rel}}$, which imprints itself on the thermodynamic history of galaxies.
% ---------------------------------------------------------
\section{Gas Accretion and the RDCC Potential}

Gas accretion onto halos is driven by the gravitational potential. In RDCC, this 
potential is modified by the relaxation mechanism:
\begin{equation}
    \Phi_{\mathrm{RDCC}}(k,a)
    = \Phi(k,a)
      \left[1 - \chi_{\Phi}
      \left(\frac{k}{k_{\mathrm{rel}}}\right)^{\alpha}\right].
\end{equation}

The suppression of small-scale modes reduces the depth of the potential wells of 
low-mass halos. As a result, the infalling gas experiences a weaker compression, 
leading to lower virial temperatures. The virial temperature in RDCC becomes
\begin{equation}
    T_{\mathrm{vir,RDCC}}
    = T_{\mathrm{vir}}
      \left[1 - \chi_{T}
      \left(\frac{k(M)}{k_{\mathrm{rel}}}\right)^{\alpha}\right].
\end{equation}

This reduction in virial temperature has profound consequences for the cooling 
efficiency of gas in low-mass halos.
% ---------------------------------------------------------
\section{Cooling Pathways and Thermodynamic Evolution}

Radiative cooling depends on the density and temperature of the gas. In RDCC, the 
modified virial temperature alters the cooling time:
\begin{equation}
    t_{\mathrm{cool,RDCC}}
    = \frac{3 k_{B} T_{\mathrm{vir,RDCC}}}{n_{e} \Lambda(T_{\mathrm{vir,RDCC}})},
\end{equation}
where $\Lambda$ is the cooling function.

Because $T_{\mathrm{vir,RDCC}}$ is reduced for low-mass halos, the cooling time 
increases, delaying the onset of star formation. Massive halos, whose collapse is 
accelerated by the coherence of large-scale modes, maintain higher virial 
temperatures and shorter cooling times. This asymmetry naturally produces a 
suppressed star formation rate in dwarf galaxies and an enhanced early star 
formation rate in massive galaxies.

The thermodynamic evolution of the gas thus becomes a direct tracer of the 
relaxation scale.
% ---------------------------------------------------------
\section{Feedback and the Regulation of Star Formation}

Feedback from supernovae and active galactic nuclei injects energy into the 
interstellar medium, reheating or expelling gas. The efficiency of feedback 
depends on the depth of the gravitational potential. In RDCC, the reduced 
potential depth of low-mass halos enhances the effectiveness of feedback, while 
the deeper potentials of massive halos remain largely unaffected.

The energy required to unbind gas from a halo of mass $M$ is
\begin{equation}
    E_{\mathrm{bind,RDCC}}
    = \frac{1}{2} M_{\mathrm{gas}} v_{\mathrm{esc,RDCC}}^{2},
\end{equation}
with
\begin{equation}
    v_{\mathrm{esc,RDCC}}
    = v_{\mathrm{esc}}
      \left[1 - \chi_{v}
      \left(\frac{k(M)}{k_{\mathrm{rel}}}\right)^{\alpha}\right].
\end{equation}

This modification amplifies the impact of feedback in dwarf-scale halos, 
suppressing their star formation even further.
% ---------------------------------------------------------
\section{Equilibrium Between Inflows and Outflows}

Galaxies evolve toward a quasi-equilibrium state in which inflows, outflows, and 
star formation balance each other. In RDCC, this equilibrium is shifted by the 
relaxation-modified potential. The equilibrium condition can be written as
\begin{equation}
    \dot{M}_{\mathrm{in}}(M)
    = \dot{M}_{\star}(M)
      + \dot{M}_{\mathrm{out}}(M),
\end{equation}
where each term carries an RDCC correction.

Low-mass halos reach equilibrium at lower stellar masses because their inflow 
rates are reduced and their outflow rates are enhanced. Massive halos reach 
equilibrium earlier and at higher stellar masses, reflecting their accelerated 
collapse and more efficient cooling.

This framework provides a natural explanation for the observed bimodality in 
galaxy populations.
% ---------------------------------------------------------
\section{Conclusion}

The baryonic cycle in the Relaxation-Driven Cyclic Cosmology is shaped by the 
scale-dependent evolution of the gravitational potential. The relaxation mechanism 
modifies gas accretion, cooling, feedback, and the equilibrium between inflows and 
outflows. These effects suppress star formation in low-mass halos, enhance early 
star formation in massive halos, and reshape the thermodynamic pathways through 
which galaxies assemble their stellar mass. The present Companion establishes the 
foundation for future studies of galaxy evolution, chemical enrichment, and the 
interaction between baryons and dark matter in RDCC.

% ========================
% References
% ========================

\begin{thebibliography}{10}

\bibitem{Flagship} Lehmann, M. (2026). Relaxation-Driven Cyclic Cosmology (RDCC): A Minimal CPT-Symmetric Framework with a Single Parameter. Flagship Paper. Zenodo: \url{https://zenodo.org/records/19744958} (v43.0 or later).

\bibitem{Zenodo} The complete RDCC companion series (I--LVII) is available at the same Zenodo record above.

% Thematisch relevante Companions
\bibitem{CompXVI} Companion XVI: Boltzmann Code and Modified Hierarchy.
\bibitem{CompXXXI} Companion XXXI: RDCC Halo Model and Non-Linear Structure.
\bibitem{CompXXXII} Companion XXXII--XXXIX: Galaxy--Halo Connection, Cosmic Web, Lensing, RSD, ISW, tSZ, Rotation Curves, CGM.

% Wichtige externe Referenzen
\bibitem{Planck2020} Planck Collaboration (2020), Astron. Astrophys. 641, A6.
\bibitem{DESI2024} DESI Collaboration (2024/2025), arXiv: [aktuelle $f\sigma_8$/BAO-Paper].
\bibitem{JWST2024} Maiolino et al. (2024), Nature (JWST high-$z$ galaxies/quasars).
\bibitem{Kaiser1987} Kaiser, N. (1987), Mon. Not. R. Astron. Soc. 227, 1 (RSD).
\bibitem{Bartelmann2001} Bartelmann, M. \& Schneider, P. (2001), Phys. Rep. 340, 291 (Weak Lensing Review).
\bibitem{HuOkamoto2002} Hu, W. \& Okamoto, T. (2002), Astrophys. J. 574, 566 (CMB Lensing).
\bibitem{LewisChallinor2006} Lewis, A. \& Challinor, A. (2006), Phys. Rep. 429, 1 (CMB Lensing Review).
\bibitem{NFW1997} Navarro, J. F., Frenk, C. S. \& White, S. D. M. (1997), Astrophys. J. 490, 493 (Halo Profiles).

\end{thebibliography}

\end{document}
